Everything in this paper is generated. The analyses live in
\texttt{paper/analysis/}; they write results as JSON into
\texttt{paper/data/} and figures into \texttt{paper/figures/}, and
\texttt{make\_numbers.py} turns that JSON into the LaTeX macros the prose
uses. No number in the text was typed by hand; re-running the
analysis rewrites the paper's numbers along with its figures.

\begin{verbatim}
make -C paper analysis   # recompute everything from the corpus
make -C paper            # build paper.pdf
make -C paper check      # fail on any LaTeX warning or bad box
\end{verbatim}

Two analyses call the live services (\texttt{/bm25}, \texttt{/semantic}) and
need them running; the rest read the index files directly. The concept-to-
passage index is expensive to build and is cached in
\texttt{paper/data/concept\_support.json}, rebuilt with
\texttt{a\_support.py --force}.

The \texttt{check} target is the reason the build is worth automating. It
fails on any LaTeX warning, any overfull or underfull box, any missing
character and any unresolved reference --- the class of defect that otherwise
ships silently in a PDF.
